\documentclass{scrartcl}
\usepackage[utf8]{inputenc}
\usepackage{hyperref}
\usepackage{url}
\usepackage{natbib}
\usepackage{graphicx}
\usepackage{float}
\usepackage{listings}
\usepackage{xcolor}
\usepackage{booktabs}
\usepackage{amsmath}
\usepackage{minted}

\usepackage{cleveref} % this must be the last package to be loaded

\title{\LARGE
Assignment \#04
}
\subtitle{Concurrent and Distributed Programming\\a.y. 2025--2026}

\author{
    Buda Marco \and Merighi Daniele \and Turchi Jacopo
}

\date{\today}

\begin{document}

\maketitle

\section{Distributed Smart Home Alarm System}

\section{Distributed TTT with Java RMI}

\section{Distributed Critical Sections with a Message-Oriented Middleware}

\subsection{Problem Analysis}
The objective is to ensure mutually exclusive access to shared resources in a distributed environment without processes having explicit knowledge of one another. From a concurrency perspective, the main challenges are:
\begin{itemize}
    \item \textbf{Race Conditions:} Simultaneous lock requests from distributed nodes must be serialized to prevent multiple processes from entering a critical section for the same resource.
    \item \textbf{Deadlocks and Faults:} A process crashing or disconnecting while holding a lock can halt the system permanently.
    \item \textbf{State Synchronization:} Lacking shared memory, the system relies entirely on asynchronous message passing, requiring precise correlation between requests and replies.
    \item \textbf{Starvation:} Queued requests must be managed to ensure fair access when resources become available.
\end{itemize}

\subsection{Adopted Strategy}
The solution implements a centralized Lock Manager utilizing RabbitMQ, following an asynchronous RPC pattern.
\begin{itemize}
    \item \textbf{Architecture:} A centralized server manages lock states.
    Clients interact via a centralized request exchange.
    \item \textbf{Message Routing:} Clients publish acquire and release requests to a shared queue.
    The server replies with grant messages to client-specific private queues, matching messages via correlation IDs.
    \item \textbf{Concurrency Control:} The server utilizes a single-thread executor to consume messages.
    This guarantees sequential processing, eliminating race conditions on the internal state maps and queues.
    \item \textbf{Hierarchical Locks:} Resources are identified by path-based targets.
    The server evaluates conflicts by checking if paths are identical or if one is a direct ancestor or descendant.
    \item \textbf{Fault Tolerance:} Locks are granted with a 15-second lease duration.
    A scheduled timer evicts expired locks, mitigating deadlocks caused by unresponsive clients.
    \item \textbf{Client Synchronization:} The client leverages \texttt{CompletableFuture} paired with a concurrent map to block execution until the asynchronous grant message is delivered or a timeout occurs.
\end{itemize}

\subsection{Execution Instructions}
To deploy and run the system, execute the following steps:
\begin{itemize}
    \item \textbf{Start RabbitMQ:} Launch the message broker container using Docker:
    \begin{minted}{bash}
    docker run -it --rm --name rabbitmq \
      -p 5672:5672 -p 15672:15672 rabbitmq:4-management
    \end{minted}
    \item \textbf{Start the Server:} Execute the centralized lock server management application (\texttt{RabbitMQLockServer}) to listen for incoming requests.
    \item \textbf{Client Lock Utilization:} Distributed processes must instantiate and invoke the provided client interface (\texttt{DistributedLockManager}) to request and release locks on target resources.
\end{itemize}

\end{document}